\documentclass{article}
\usepackage[utf8]{inputenc}
\usepackage{amsmath,amssymb}
\usepackage[most]{tcolorbox}
\usepackage{geometry}
\geometry{margin=2.5cm}

\newcommand{\step}[1]{\par\noindent\hangindent=3.5em\hangafter=1%
  \makebox[3.5em][l]{\textbf{#1.}}}
\newcommand{\steptag}[1]{\hfill{\small\textsf{[#1]}}}

\newtcolorbox{proofbox}[1]{
  colback=gray!5, colframe=gray!60,
  fonttitle=\bfseries, title={#1},
  breakable, enhanced,
  left=4pt, right=4pt, top=4pt, bottom=4pt,
  before skip=10pt, after skip=10pt
}

\begin{document}

\begin{proofbox}{After first fixing one global witness package W\_RF suppli...}
\step{1} After first fixing one global witness package W\_RF supplied by lem-routef-raw-factor-setting-formation, for every input (H,Phi,eta) to which that formation result applies, fix one def-routef-raw-factor-setting datum S over that same W\_RF supplied by the same result, for every Delta' supplied for (W\_RF,S) by lem-routef-delta-prime-closeness, every Delta supplied from that same Delta' by lem-routef-delta-normalization-closeness, every Upsilon' supplied from that same pair by lem-routef-upsilon-prime-closeness, and every Upsilon supplied from that same triple by lem-routef-upsilon-normalization-closeness; for every integer n >= 1 and all X, Y in M\_n(S.B), writing the fields of (W\_RF,S) as the unqualified symbols below: Multiplicative telescope: for rho\_mult := min\{rho\_T, rho\_id, rho\_DeltaPhi, rho\_Upsilon\} and 0 <= eta <= rho\_mult, every amplification satisfies ||Upsilon\_n(Delta\_n X Delta\_n Y) - XY|| <= [C\_Upsilon+2*(C\_2+C\_theta+C\_Delta)]*eta*||X||*||Y||. \steptag{validated} \\[6pt]
\step{1.1} Radius and operator-norm consequences. Fix the globally chosen W\_RF, an admissible input and its chosen S, Delta', Delta, Upsilon', Upsilon, and fix n>=1 and X,Y in M\_n(B). Assume 0<=eta<=rho\_mult. By def-routef-raw-factor-setting (1.1),(1.2),(1.5), rho\_mult=min\{rho\_T,rho\_id,rho\_DeltaPhi,rho\_Upsilon\}, rho\_DeltaPhi=min\{rho\_theta,rho\_Delta,rho\_2\}, rho\_id\textasciicircum{}corr=min\{rho\_theta,rho\_AI,epsilon\_E/C\_A\}, rho\_id=min\{rho\_AI,epsilon\_E/C\_A\}, rho\_T<=rho\_theta, and C\_T=C\_theta+3*C\_V. Hence eta is in the ranges of lem-routef-raw-factor-norms, lem-routef-raw-factor-identities, lem-routef-delta-normalization-closeness, lem-routef-delta-phi-product, and lem-routef-upsilon-normalization-closeness. The first gives ||tilde-Upsilon||\_cb<=1+C\_T*eta. Moreover rho\_T<=1/[4(1+C\_theta)] and rho\_T<=1/[4(1+C\_V)], while C\_theta,C\_V>=0, so C\_T*eta<=C\_theta/[4(1+C\_theta)]+3*C\_V/[4(1+C\_V)]<=1 and therefore ||tilde-Upsilon||\_cb<=2. The normalization lemma says Delta is UCP; consequently Delta is completely contractive: for each amplification, unital complete positivity and the 2-by-2/Schwarz argument give Delta\_n(Z)\textasciicircum{}*Delta\_n(Z)<=Delta\_n(Z\textasciicircum{}*Z)<=||Z||\textasciicircum{}2 I, hence ||Delta\_n(Z)||<=||Z||. In particular ||Delta\_n X Delta\_n Y||<=||X||*||Y||. \steptag{validated} \\[4pt]
\step{1.2} Exact telescope identity. Under the same choices and radius, put Z:=Delta\_n(X)Delta\_n(Y) and E:=tilde-Phi\_n(Z)-tilde-Delta\_n(XY). By lem-routef-raw-factor-identities, tilde-Delta tilde-Upsilon=tilde-Phi and tilde-Upsilon tilde-Delta=I\_B; amplification preserves these composition identities. Thus tilde-Upsilon\_n tilde-Phi\_n=tilde-Upsilon\_n tilde-Delta\_n tilde-Upsilon\_n=tilde-Upsilon\_n and tilde-Upsilon\_n tilde-Delta\_n(XY)=XY. Adding and subtracting tilde-Upsilon\_n(Z) therefore gives the exact equality Upsilon\_n(Z)-XY=(Upsilon\_n-tilde-Upsilon\_n)(Z)+tilde-Upsilon\_n(E). \steptag{validated} \\[4pt]
\step{1.3} Estimate the two telescope terms and conclude. By lem-routef-upsilon-normalization-closeness, ||Upsilon-tilde-Upsilon||\_cb<=C\_Upsilon*eta, so complete-contractivity of Delta yields ||(Upsilon\_n-tilde-Upsilon\_n)(Z)||<=C\_Upsilon*eta*||X||*||Y||. By lem-routef-delta-phi-product, ||E||<=(C\_2+C\_theta+C\_Delta)*eta*||X||*||Y||, and the bound ||tilde-Upsilon||\_cb<=2 yields ||tilde-Upsilon\_n(E)||<=2*(C\_2+C\_theta+C\_Delta)*eta*||X||*||Y||. Applying the triangle inequality to the exact identity gives ||Upsilon\_n(Delta\_n X Delta\_n Y)-XY||<=[C\_Upsilon+2*(C\_2+C\_theta+C\_Delta)]*eta*||X||*||Y||, as required for every n,X,Y. \steptag{validated} \\[6pt]
\step{1.3.1} Local radius consequences and first-term estimate. Under the hypotheses and notation of root node 1, fix n>=1 and X,Y in M\_n(B), assume 0<=eta<=rho\_mult, and define Z:=Delta\_n(X)Delta\_n(Y). From rho\_mult=min\{rho\_T,rho\_id,rho\_DeltaPhi,rho\_Upsilon\}, rho\_DeltaPhi=min\{rho\_theta,rho\_Delta,rho\_2\}, rho\_id=min\{rho\_AI,epsilon\_E/C\_A\}, and rho\_id\textasciicircum{}corr=min\{rho\_theta,rho\_AI,epsilon\_E/C\_A\}, we have eta<=rho\_T, eta<=rho\_Delta, eta<=rho\_DeltaPhi, eta<=rho\_Upsilon, and, because rho\_T<=rho\_theta, eta<=rho\_id\textasciicircum{}corr. Thus the stated upstream lemmas apply. In particular, lem-routef-delta-normalization-closeness makes Delta UCP, so Delta\_n is unital completely positive and hence satisfies Kadison-Schwarz: Delta\_n(T)\textasciicircum{}*Delta\_n(T)<=Delta\_n(T\textasciicircum{}*T)<=||T||\textasciicircum{}2 I. Therefore each Delta\_n is contractive and ||Z||<=||X||*||Y||. Also lem-routef-raw-factor-norms gives ||tilde-Upsilon||\_cb<=1+C\_T*eta. Here C\_theta=12*(sqrt(2)-1)>0, bar C\_E=max\{1,C\_E\}>=1, C\_A=20+(211/8)*C\_theta>0, and C\_V=bar C\_E*C\_A>=0. Since C\_T=C\_theta+3*C\_V, eta<=rho\_T<=1/[4(1+C\_theta)] and eta<=rho\_T<=1/[4(1+C\_V)], we get C\_T*eta<=C\_theta/[4(1+C\_theta)]+3*C\_V/[4(1+C\_V)]<=1, hence ||tilde-Upsilon||\_cb<=2. Finally lem-routef-upsilon-normalization-closeness gives ||Upsilon-tilde-Upsilon||\_cb<=C\_Upsilon*eta, so ||(Upsilon\_n-tilde-Upsilon\_n)(Z)||<=C\_Upsilon*eta*||X||*||Y||. \steptag{validated} \\[4pt]
\step{1.3.2} Local error definition, exact identity, and second-term estimate. Under the hypotheses of root node 1 and with Z from node 1.3.1, define E:=tilde-Phi\_n(Z)-tilde-Delta\_n(XY). The radius conclusions in node 1.3.1 permit lem-routef-raw-factor-identities and lem-routef-delta-phi-product. Amplifying tilde-Delta tilde-Upsilon=tilde-Phi and tilde-Upsilon tilde-Delta=I\_B gives tilde-Upsilon\_n tilde-Phi\_n=tilde-Upsilon\_n tilde-Delta\_n tilde-Upsilon\_n=tilde-Upsilon\_n and tilde-Upsilon\_n tilde-Delta\_n(XY)=XY. Hence, by direct expansion, Upsilon\_n(Z)-XY=(Upsilon\_n-tilde-Upsilon\_n)(Z)+tilde-Upsilon\_n(tilde-Phi\_n(Z)-tilde-Delta\_n(XY))=(Upsilon\_n-tilde-Upsilon\_n)(Z)+tilde-Upsilon\_n(E). Moreover lem-routef-delta-phi-product gives ||E||<=(C\_2+C\_theta+C\_Delta)*eta*||X||*||Y||, and node 1.3.1 gives ||tilde-Upsilon||\_cb<=2; therefore ||tilde-Upsilon\_n(E)||<=2*(C\_2+C\_theta+C\_Delta)*eta*||X||*||Y||. \steptag{validated} \\[4pt]
\step{1.3.3} Conclusion from the two explicit telescope bounds. By the exact equality in node 1.3.2, the triangle inequality, the first-term estimate in node 1.3.1, and the second-term estimate in node 1.3.2, ||Upsilon\_n(Delta\_n X Delta\_n Y)-XY||=||Upsilon\_n(Z)-XY||<=[C\_Upsilon+2*(C\_2+C\_theta+C\_Delta)]*eta*||X||*||Y||. Since n>=1 and X,Y were arbitrary under the root hypotheses, this is precisely the conclusion asserted by node 1.3. \steptag{validated} \\[4pt]
\end{proofbox}

\end{document}
